% chapters/12/section-05.tex
\section{Skrip C\# Pertama}
\label{sec:skrip-pertama}

\subsection{Kode Lengkap}

\kodecsharp{unity/ch12_skrip_pertama.cs}{Skrip C\# pertama: menggerakkan, memutar, dan melaporkan keadaan benda}{lst:ch12-skrip}

\subsection{Menempelkan Skrip pada Benda}

Skrip yang baru ditulis belum berpengaruh apa pun. Ia perlu ditempelkan pada
sebuah \istilah{GameObject} lebih dahulu, dan caranya ada dua.

\begin{enumerate}
  \item \textbf{Menyeret berkas skrip} dari \istilah{Project window} ke sebuah
        benda pada \istilah{Hierarchy} atau pada \istilah{Scene}.
  \item \textbf{Menambahkan komponen} melalui tombol \istilah{Add Component}
        pada \istilah{Inspector}, lalu memilih nama skripnya.
\end{enumerate}

Setelah ditempelkan, skrip itu tampak sebagai salah satu komponen pada
\istilah{Inspector}, beserta nilai pengaturannya. Nilai itulah yang Anda ubah
untuk mengatur kecepatan geraknya, tanpa menyunting kode.

\subsection{Fungsi yang Dipanggil Mesin}

Skrip tersebut tidak memuat siklus permainan. Tidak ada
perulangan \code{while}, dan tidak ada pemanggilan gambar. Yang ada hanyalah
dua fungsi, yaitu \code{Start} dan \code{Update}, beserta fungsi ketiga yaitu
\code{OnDisable}.

Fungsi-fungsi itu \textbf{dipanggil oleh mesin} pada saat yang telah
ditentukan. Urutannya tetap dan sama untuk setiap benda, sebagaimana
diperlihatkan pada Tabel~\ref{tab-daur-hidup}.

\begin{table}[htbp]
\centering
\caption{Fungsi skrip yang dipanggil mesin beserta waktu pemanggilannya}
\label{tab-daur-hidup}
{\setlength{\tabcolsep}{4pt}\footnotesize%
\begin{tabularx}{\textwidth}{@{}L{3.0cm} L{3.4cm} Y@{}}
\toprule
\headrow \textbf{Fungsi} & \textbf{Kapan dipanggil} & \textbf{Dipakai untuk} \\
\midrule
\code{Awake} & Sekali, ketika benda dibuat &
Menyiapkan rujukan ke komponen lain, termasuk yang tidak bergantung pada benda
lain. \\
\addlinespace
\code{Start} & Sekali, sebelum bingkai pertama &
Menyiapkan nilai awal yang bergantung pada benda lain, yang belum tentu siap
ketika \code{Awake} dijalankan. \\
\addlinespace
\code{Update} & Sekali setiap bingkai &
Menangani masukan dan hal yang mengikuti lama bingkai. Padanan tahap pembaruan
pada Bab~IV. \\
\addlinespace
\code{FixedUpdate} & Berkala, dengan selang tetap &
Menangani perhitungan fisika. Padanan \istilah{fixed timestep} pada Bagian~2.2. \\
\addlinespace
\code{OnDisable} & Ketika skrip dimatikan &
Membereskan pekerjaan, misalnya menyimpan keadaan atau menutup berkas. \\
\bottomrule
\end{tabularx}}
\end{table}

Urutan pemanggilan itu menjawab pertanyaan yang telah Anda
kenal sejak Bab~IV: mana yang lebih dahulu, membaca masukan atau memperbarui
keadaan? Pada Unity, mesin yang menentukan urutannya, sedangkan tugas Anda
adalah meletakkan pekerjaan pada fungsi yang tepat.

\begin{catatabox}
Perhatikan pemakaian \code{Time.deltaTime} pada skrip di atas. Nilainya
menyatakan lama bingkai yang baru saja berlangsung, sebagaimana
\code{delta_seconds} pada Bab~V. Perhatikan pula pemakaian
\code{direction.normalized}: penormalan vektor arah yang Anda pelajari pada
Bagian~5.2 tetap diperlukan, sebab tanpa itu benda bergerak lebih cepat ketika
menempuh arah diagonal.
\end{catatabox}
